\documentclass[11pt,a4paper]{article}
\usepackage[margin=1.7cm]{geometry}
\usepackage{amsmath,amssymb,mathtools}
\usepackage{xcolor}
\usepackage{tcolorbox}
\usepackage{enumitem}
\usepackage{booktabs}
\usepackage{array}
\usepackage{hyperref}
\usepackage{fancyhdr}
\usepackage{lmodern}
\usepackage[T1]{fontenc}
\hypersetup{colorlinks=true,linkcolor=blue!50!black,urlcolor=blue!50!black}
\pagestyle{fancy}
\fancyhf{}
\lhead{MATH6342 Linear Systems}
\rhead{Week 1}
\cfoot{\thepage}
\setlength{\parindent}{0pt}
\setlength{\parskip}{5pt}
\definecolor{navy}{HTML}{243B6B}
\definecolor{lightblue}{HTML}{EEF4FF}
\definecolor{lightgreen}{HTML}{EEF9F0}
\definecolor{lightyellow}{HTML}{FFF8E6}
\definecolor{lightred}{HTML}{FFF0F0}
\newtcolorbox{keybox}[1][]{colback=lightblue,colframe=navy,title=#1,boxrule=0.8pt,arc=2mm}
\newtcolorbox{goalbox}[1][]{colback=lightgreen,colframe=green!45!black,title=#1,boxrule=0.8pt,arc=2mm}
\newtcolorbox{warnbox}[1][]{colback=lightyellow,colframe=orange!70!black,title=#1,boxrule=0.8pt,arc=2mm}
\newtcolorbox{examplebox}[1][]{colback=lightred,colframe=red!55!black,title=#1,boxrule=0.8pt,arc=2mm}
\newcommand{\vect}[1]{\boldsymbol{#1}}

\begin{document}

\begin{center}
{\Huge\bfseries MATH6342 Linear Systems}\\[4pt]
{\LARGE Week 1: Two 1-Hour Lectures}\\[8pt]
{\large Simplified preview based on the opening Chapter 1 material in the supplied lecture notes}\\[10pt]
\end{center}

\begin{keybox}[How this guide was estimated]
The source notes are organised by \textbf{frames/topics}, not by teaching week. Since you said there are \textbf{two one-hour lectures per week}, this guide treats the opening chapter, \emph{Linear systems and their description} (frames 6--30), as the most likely Week 1 block. The exact stopping point may vary slightly in class.
\end{keybox}

\section*{The big picture for Week 1}
The first week is likely about learning how a physical or mathematical process can be represented as a \textbf{linear system}. The basic language is
\[
\boxed{\text{input }u(t)}\;\longrightarrow\;\boxed{\text{system}}\;\longrightarrow\;\boxed{\text{output }y(t)}.
\]
A system may also have an internal \textbf{state} $x(t)$, which stores enough information about the system at time $t$ to determine how it evolves in the future.

The most important progression is likely to be
\[
\text{physical system}
\longrightarrow \text{block diagram}
\longrightarrow \text{differential equation}
\longrightarrow \boxed{X'=AX+Bu}.
\]

\section*{Lecture 1 (about 1 hour): What is a linear system?}

\begin{goalbox}[Likely learning goals]
By the end of Lecture 1 you should be able to identify the input, output and state of a system, recognise basic linear operators, read simple block diagrams, and distinguish open-loop from feedback systems.
\end{goalbox}

\subsection*{1. Input, output and state}
The notes describe a system evolving in time $t>0$ using:
\begin{itemize}[leftmargin=1.5em]
\item \textbf{Input} $u(t)$: what is supplied to or drives the system.
\item \textbf{Output} $y(t)$: what we observe or measure.
\item \textbf{State} $x(t)$: the internal variables needed to determine future behaviour.
\end{itemize}

A key question throughout the module is:
\[
\boxed{\text{How does }y(t)\text{ depend on }u(t)?}
\]

\subsection*{2. Linear operators}
You will revisit the definition of a linear map:
\[
L(\lambda f+\mu g)=\lambda Lf+\mu Lg.
\]
The opening examples include:
\begin{center}
\begin{tabular}{>{\bfseries}p{4cm}p{8.5cm}}
\toprule
Operator & Effect \\
\midrule
Differentiator & $f\mapsto \dfrac{df}{dt}$ \\
Integrator & $f\mapsto \displaystyle\int_0^x f(t)\,dt$ \\
Amplifier & $f\mapsto af$ \\
Multiplication & $f(t)\mapsto h(t)f(t)$ \\
Evaluation & $f\mapsto f(t_0)$ \\
\bottomrule
\end{tabular}
\end{center}

\subsection*{3. Block-diagram building blocks}
You will likely learn how to represent systems visually using:
\begin{itemize}[leftmargin=1.5em]
\item \textbf{Summing junction}: combines signals, e.g. $u_1+u_2$.
\item \textbf{Tap}: copies the same signal into more than one branch.
\item \textbf{Amplifier}: $u\mapsto au$.
\item \textbf{Integrator}: $u\mapsto \int u$.
\item \textbf{Differentiator}: $u\mapsto \dfrac{du}{dt}$.
\end{itemize}

\begin{examplebox}[Typical block-diagram equation]
The notes give the example
\[
y=3\int u+2u'+6u.
\]
You should be able to read this as three parallel branches acting on the same input $u$, followed by a summing junction.
\end{examplebox}

\subsection*{4. Open-loop and feedback systems}
An \textbf{open-loop} system processes the input and produces the output without sending the output back into the system.

A \textbf{feedback} (closed-loop) system sends some of the output back through the system. This is central to control theory.

\begin{center}
\begin{tabular}{p{6cm}p{6cm}}
\toprule
\textbf{Positive feedback} & \textbf{Negative feedback} \\
\midrule
Reinforces changes and can amplify growth. & Opposes changes and can stabilise behaviour. \\
Bird population: $\dfrac{dx}{dt}=kx+u$. & Cyclist with air resistance: $m\dfrac{dv}{dt}=-kv+u$. \\
\bottomrule
\end{tabular}
\end{center}

\begin{warnbox}[What matters most in Lecture 1]
Do not just memorise the symbols. Practise translating in both directions:
\[
\text{equation} \leftrightarrow \text{block diagram}.
\]
That skill is likely to make the later state-space material much easier.
\end{warnbox}

\newpage
\section*{Lecture 2 (about 1 hour): From differential equations to state-space form}

\begin{goalbox}[Likely learning goals]
By the end of Lecture 2 you should understand how linear differential equations can be represented as systems, why matrix notation is useful, and how a higher-order scalar ODE can be rewritten as a first-order vector equation.
\end{goalbox}

\subsection*{1. Matrix amplifiers and vector signals}
The course extends scalar block diagrams to vector inputs. A matrix block acts as
\[
U\mapsto BU.
\]
This allows the same diagram language to handle \textbf{multiple inputs and multiple outputs}.

\subsection*{2. Linear differential operators}
The notes introduce operators of the form
\[
Lf(t)=a_n(t)\frac{d^nf}{dt^n}+a_{n-1}(t)\frac{d^{n-1}f}{dt^{n-1}}+\cdots+a_0(t)f(t).
\]
If the coefficients $a_j$ are constants, this is a \textbf{constant-coefficient linear differential operator}.

A typical equation is
\[
a_n y^{(n)}+a_{n-1}y^{(n-1)}+\cdots+a_0y
=
 b_m u^{(m)}+\cdots+b_0u.
\]
The notes show that such an equation can be realised as a feedback system using integrators, differentiators, amplifiers, taps and summing junctions.

\subsection*{3. The damped harmonic oscillator}
This is one of the key examples and is likely to reappear throughout the module:
\[
\boxed{y''+\beta y'+\gamma y=u.}
\]
Interpretation:
\begin{itemize}[leftmargin=1.5em]
\item $y$ = displacement,
\item $y'$ = velocity,
\item $y''$ = acceleration,
\item $u(t)$ = external driving force,
\item $-\beta y'$ = damping,
\item $-\gamma y$ = restoring force.
\end{itemize}
The notes use car suspension as the motivating physical example.

\subsection*{4. The key conversion: second-order ODE to first-order system}
Set
\[
v=y'.
\]
Then
\[
\begin{aligned}
y'&=v,\\
v'&=-\gamma y-\beta v+u.
\end{aligned}
\]
Now define the state vector
\[
X=\begin{pmatrix}y\\v\end{pmatrix}.
\]
Then the two equations combine into
\[
\boxed{\frac{dX}{dt}=AX+Bu},
\]
with
\[
A=\begin{pmatrix}
0&1\\
-\gamma&-\beta
\end{pmatrix},
\qquad
B=\begin{pmatrix}0\\1\end{pmatrix}.
\]

\begin{keybox}[Why this is such a big deal]
A second-order scalar equation has been converted into a \textbf{first-order matrix system}. This means the tools of linear algebra can now be used to study differential equations and dynamical systems.
\end{keybox}

\subsection*{5. General reduction of order}
The same idea works for an $n$th-order equation. Define
\[
X=\begin{pmatrix}
y\\y'\\\vdots\\y^{(n-1)}
\end{pmatrix}.
\]
Then the $n$th-order scalar ODE can be written as
\[
\boxed{X'=AX+Bu}.
\]
The matrix $A$ has the familiar shift structure with 1s just above the diagonal and the differential-equation coefficients in the last row. This leads naturally to the \textbf{companion matrix}, which appears at the end of the opening chapter.

\section*{What you should know by the end of Week 1}
You should be comfortable with the following checklist:

\begin{enumerate}[leftmargin=1.8em]
\item Explain what \textbf{input}, \textbf{output} and \textbf{state} mean.
\item State the definition of a \textbf{linear map}.
\item Recognise an amplifier, integrator, differentiator, tap and summing junction.
\item Translate a simple equation into a block diagram and back again.
\item Explain the difference between \textbf{open-loop} and \textbf{feedback} systems.
\item Distinguish positive from negative feedback in simple examples.
\item Recognise a linear differential operator.
\item Rewrite a second-order ODE as two first-order equations.
\item Form the state vector $X$ and identify the matrices $A$ and $B$.
\item Understand the central state-space form
\[
\boxed{X'=AX+Bu}.
\]
\end{enumerate}

\section*{Best things to revise before Week 1}
You do not need to pre-learn control theory. The most useful background is:
\begin{itemize}[leftmargin=1.5em]
\item matrix multiplication and matrix-vector multiplication;
\item $2\times2$ matrices;
\item basic differentiation and integration;
\item first- and second-order ODE notation;
\item the definition of a linear map;
\item changing a second-order ODE into a pair of first-order equations.
\end{itemize}

\section*{Quick self-test}
Try these without looking back:
\begin{enumerate}[leftmargin=1.8em]
\item What is the difference between an input $u(t)$ and a state $x(t)$?
\item What does an amplifier labelled $a$ do to a signal $u$?
\item What is the difference between a tap and a summing junction?
\item Why is $m v'=-kv+u$ an example of negative feedback?
\item For $y''+4y'+3y=u$, let $v=y'$. Write down the two first-order equations.
\item For the same equation, write down $A$ and $B$ in $X'=AX+Bu$.
\end{enumerate}

\begin{warnbox}[Answer to Questions 5--6]
For $y''+4y'+3y=u$,
\[
y'=v,\qquad v'=-3y-4v+u,
\]
so
\[
X=\begin{pmatrix}y\\v\end{pmatrix},\qquad
A=\begin{pmatrix}0&1\\-3&-4\end{pmatrix},\qquad
B=\begin{pmatrix}0\\1\end{pmatrix}.
\]
\end{warnbox}

\vfill
\begin{center}
\small Prepared as a Week 1 preview from the supplied MATH6342 lecture notes and beamer slides.\\
The source materials are organised by frame rather than by teaching week, so the two-lecture split is an informed estimate.
\end{center}

\end{document}
